% GENERATED, edit the French unit or the translation table
\mxunit{terminale-ateliers}{End-of-year and oral workshops}
\section{O1 · Turning a theme into a question}
\textbf{S82 · 2 h.} Choose a situation that allows reasoning from the syllabus, then write a precise question. Examples: ``How do we choose a number of trials to make a failure rare?''; ``How can a cooling model be used to predict a waiting time?''; ``Why does a mean become more stable as the sample grows?''
Produce a page containing the question, the hypotheses, a numerical example, an accessible proof and a limitation of the model. Allow 20 min of individual research, 30 min of mathematical work, 30 min of group discussion, 25 min of writing and 15 min of feedback.
\section{O2 · Presenting and answering}
\textbf{S84 · 2 h.} Prepare a 10 min presentation and answers to the following questions: why this model? which hypothesis is decisive? what does each variable mean? which result is proved? what changes with another value of the parameter?
Practise in groups of three: a speaker, an interviewer and an observer. The observer notes one clear step, one step to justify and one wording to improve. Swap roles. A video is useful only if it is permitted and actually watched back.
\section{O3 · Oral simulation}
\textbf{S87 and S93.} Prepare two questions and give a presentation followed by a discussion. Use a simple grid: mathematical accuracy, structure, explanation of notation, response to objections, quality of the exchange. The final arrangements are those of the instructions and summons for the session.
\ifmxteacher\begin{teachernote}[Organising the sittings]
A 20 min sitting takes a third of an hour: at most six sittings in 2 h, not counting transitions. With a whole class, plan several panels or spread out the sittings; do not announce a full individual oral for everyone in a single slot. Have students specify what is proved and what is assumed in each model. An internal grid is used for teaching feedback; it does not replace the official grid for the session.
\end{teachernote}\fi
